\section*{第 \ref{ch:b3:probability} 章——概率：
基础与大数定律}

\begin{solution}{exo:b3:probability:1}
(a) 对 $u \in \intoo01$，$\P(F(X) \leq u) = \P(X \leq
F^{-1}(u)) = F(F^{-1}(u)) = u$（连续性与严格单调性使 $F$
成为 $\intoo01$ 上的双射，且
$\{F(X) \leq u\} = \{X \leq F^{-1}(u)\}$）：$F(X)$ 服从均匀分布。
反之，$\P(G(U) \leq t) = \P(U \leq F(t)) = F(t)$：欲模拟某一分布，
只需将逆分布函数作用于一个
均匀样本。

(b) 设 $Y = X^2$，$X$ 在 $\intcc{-1}1$ 上服从均匀分布：对 $t \in
\intcc01$，$F_Y(t) = \P(-\sqrt t \leq X \leq \sqrt t) = \sqrt
t$，故密度为 $\frac1{2\sqrt t}\mathbf 1_{\intoo01}$。又
$\P\bigl(-\frac1\lambda\ln U \leq t\bigr) = \P(U \geq
\eu^{-\lambda t}) = 1 - \eu^{-\lambda t}$，即指数分布
$\mathcal E(\lambda)$——逆变换法的一个实例。
\end{solution}

\begin{solution}{exo:b3:probability:2}
(a) 泊松分布：$\E X = \sum_{k\geq0}k\,\eu^{-\lambda}
\frac{\lambda^k}{k!} = \lambda$，$\E[X(X-1)] = \lambda^2$，于是
$\V = \lambda^2 + \lambda - \lambda^2 = \lambda$。几何分布
($\P(X = k) = p(1-p)^{k-1}$)：$\E X = \frac1p$，$\V =
\frac{1-p}{p^2}$（对几何级数逐项求导两次）。

(b) $G(t) = \P(T > t)$ 单调不增，且 $G(0^+)\dots$
$G \colon \intco0\infty \to \intoc01$；无记忆性表述为
$G(t + s) = G(t)G(s)$。于是 $G(n t) = G(t)^n$，$G(t/n) =
G(t)^{1/n}$，故 $G(q) = G(1)^q$（$q \geq 0$ 为有理数）；记
$G(1) = \eu^{-\lambda}$（$\in \intoo01$：若 $G(1) = 1$，则
必有 $G \equiv 1$，对有限\omterm{def:b3:probability:space}{随机变量}不可能；
$G(1) = 0$ 与假设矛盾），再以有理数夹逼
任一 $t$（由单调性）：$G(t) =
\eu^{-\lambda t}$——即指数分布。其逆命题
只是一个计算。
\end{solution}

\begin{solution}{exo:b3:probability:3}
(a) $\sigma(f_i(X_i)) = f_i(X_i)^{-1}(\mathcal B) \subseteq
X_i^{-1}(\mathcal B) = \sigma(X_i)$（$f_i$ 为博雷尔函数），且
\omterm{def:b3:probability:independence}{独立} $\sigma$-代数的子 $\sigma$-代数
仍然\omterm{def:b3:probability:independence}{独立}（定义中的等式显然成立）。

(b) $\sigma(A_i) = \{\varnothing, A_i, A_i^c, \Omega\} =
\sigma(A_i^c) = \sigma(\mathbf 1_{A_i})$：三个陈述
断言的是同一族
$\sigma$-代数的\omterm{def:b3:probability:independence}{独立}性。（由 $A_i$ 上的这一分解推广到补集，
正是 $\lambda$-系论证，它位于
\cref{def:b3:probability:independence} 的等价性之中
——或直接用容斥公式。）

(c) $\P(A) = \P(B) = \P(C) = \frac12$；$A\cap B = A\cap C =
B\cap C$ 对各对而言：每个交事件为“两次均为正面”或
类似事件，概率均为 $\frac14$，故两两\omterm{def:b3:probability:independence}{独立}。
但 $\P(A\cap B\cap C) = \P(\text{HH}) = \frac14 \neq
\frac18$：不\omterm{def:b3:probability:independence}{独立}——$C$ 由 $A$ 与
$B$ 决定。
\end{solution}

\begin{solution}{exo:b3:probability:4}
(a) 设文本 $T$ 的长度为 $L$，$q = a^{-L}$（$a$
为字母表大小）。事件 $E_k = \{$位置 $kL+1,
\dots, (k+1)L$ 拼出 $T\}$ 相互\omterm{def:b3:probability:independence}{独立}（不相交的块
由\omterm{def:b3:probability:independence}{独立}同分布的字母组成），每个概率为 $q > 0$：$\sum\P(E_k)
= \infty$，由博雷尔–坎泰利引理（2）得
几乎必然出现无穷多次。

(b) 上界：$\P\bigl(R_n \geq (1+\varepsilon)\log_2n\bigr)
\leq 2^{-(1+\varepsilon)\log_2n} = n^{-(1+\varepsilon)}$，
可和；由博雷尔–坎泰利引理（1），几乎必然只有有限个
这样的 $n$。下界：取互不相交的块——第 $j$ 个块
长度为 $\ell_j = \lceil\log_2s_j\rceil$，起点为 $s_j =
\sum_{i<j}\ell_i$。事件“第 $j$ 块全为 1”
相互\omterm{def:b3:probability:independence}{独立}，概率为 $2^{-\ell_j} \asymp
\frac1{s_j} \asymp \frac1{j\log_2 j}$，其和发散；
由博雷尔–坎泰利引理（2）得全 1 块出现无穷多次，
即 $R_{s_j} \geq \log_2 s_j$ 无穷多次成立。综上，
前 $n$ 位数字中的最长游程长度
几乎必然为 $(1 + o(1))\log_2n$。
\end{solution}

\begin{solution}{exo:b3:probability:5}
(a) $R = \bigl(\limsup\abs{X_n}^{1/n}\bigr)^{-1}$ 在
有限个 $X_n$ 被修改时保持不变：对每个 $N$，
$R$ 关于 $\sigma(X_N, X_{N+1}, \dots)$ 可测，即
尾可测。于是每个事件 $\{R \leq c\}$
的概率为 $0$ 或 $1$（由\cref{thm:b3:probability:zeroone}），
故 $R$ 的分布函数只取
$0, 1$ 两个值：它恰在一点 $c_0 \in
\intcc0{+\infty}$ 处跳跃，从而 $R = c_0$ 几乎必然成立。

(b) $\sum X_n$ 的收敛性与 $\frac{S_n}n$ 的收敛性
对有限项的改动不敏感（对后者而言：
被改动的项贡献 $O(1/n) \to 0$）。二者均为尾事件，
适用 0–1 律。

(c) $\{X_1 > 0\}$ 依赖于 $X_1$：对\omterm{def:b3:probability:independence}{独立}同分布的符号
（$\P(X_1 = \pm1) = \frac12$），其概率为 $\frac12
\notin \{0,1\}$——并无矛盾，它不是
尾事件。
\end{solution}

\begin{solution}{exo:b3:probability:6}
在 $(\intcc01, \lambda)$ 上进行讨论。
(a) 打字机序列 $\mathbf 1_{I_n}$
（见\cref{exo:b3:lp:3}）：$\norm{X_n}_p^p = \lambda(I_n) \to 0$
（对所有 $p < \infty$），因此也依概率收敛；在每个
$\omega$ 处，取值 $0$ 与 $1$ 都会反复出现：处处不存在
逐点收敛。
(b) $X_n = n\mathbf 1_{\intoo0{1/n}} \to 0$（除 $0$ 外），但
$\norm{X_n}_p \geq n^{1 - 1/p} \geq 1$。
(c) $X_n = \sqrt n\,\mathbf 1_{\intoo0{1/n}}$：$\E\abs{X_n} =
n^{-1/2} \to 0$，$\E X_n^2 = 1$。
(d) 对任一子列，依概率收敛可再抽取一个
几乎必然收敛的子列（由\cref{prop:b3:probability:modes}(c)）；由控制收敛定理
沿该子列得到 $L^1$ 收敛，
且极限仍是\emph{同一个}
$X$。于是数值数列 $\E\abs{X_n - X}$ 的每个子列
都有趋于 $0$ 的子子列：整个数列
收敛到 $0$。
\end{solution}

\begin{solution}{exo:b3:probability:7}
(a) $\hat p_n = \frac{S_n}n$，其中 $S_n$ 服从二项分布，$\V(\hat
p_n) = \frac{p(1-p)}n \leq \frac1{4n}$，由切比雪夫
（\cref{prop:b3:probability:markov}）即得所求界。
(b) 解 $\frac1{4n(0.03)^2} \leq 0.05$：$n \geq
\frac{1}{4\cdot0.0009\cdot0.05} \approx 5556$。中心
极限定理可说明，同样的保证只需 $n \approx 1070$，
而切比雪夫以一般性为代价，其因子
约为 $\approx 5$。
\end{solution}

\begin{solution}{exo:b3:probability:8}
(a) 若 $S_n \sim \mathcal B(n, x)$，由变换定理
可得 $\E\bigl[f(\frac{S_n}n)\bigr] =
\sum_k\binom nkx^k(1-x)^{n-k}f(\frac kn) = B_nf(x)$。

(b)--(c) 设 $\omega = \omega_f$ 为连续模
（$\abs{f(u) - f(v)} \leq \omega(\abs{u - v})$，且 $\omega(c
\delta) \leq (1 + c)\,\omega(\delta)$，由分段拼接可得）。
于是对任意 $\delta > 0$，
\[
\abs{f(u) - f(x)} \leq \Bigl(1 + \frac{(u -
x)^2}{\delta^2}\Bigr)\omega(\delta)
\]
（若 $\abs{u - x} \leq \delta$，显然成立；否则
$\omega(\abs{u-x}) \leq (1 + \frac{\abs{u-x}}\delta)
\omega(\delta) \leq (1 + \frac{(u-x)^2}{\delta^2})
\omega(\delta)$）。在 $u = \frac{S_n}n$ 处取\omterm{def:b3:probability:space}{期望}：
\[
\abs{B_nf(x) - f(x)} \leq
\Bigl(1 + \frac{\V(S_n/n)}{\delta^2}\Bigr)\omega(\delta)
\leq \Bigl(1 + \frac{1}{4n\delta^2}\Bigr)\omega(\delta) ;
\]
取 $\delta = n^{-1/2}$，则 $\norm{B_nf - f}_\infty \leq
\frac54\,\omega\bigl(n^{-1/2}\bigr) \leq
\frac32\,\omega\bigl(n^{-1/2}\bigr) \to 0$（一致
连续性，紧集上）：概率形式的魏尔斯特拉斯
定理，并给出显式且一致的速率。
\end{solution}

\begin{solution}{exo:b3:probability:9}
(a) 收集到 $k - 1$ 种类型后，每次抽取为新类型的概率为
$p_k = \frac{n-k+1}n$：$\tau_k$ 服从几何分布，
参数为 $(p_k)$，且各 $\tau_k$ 相互\omterm{def:b3:probability:independence}{独立}（因抽取\omterm{def:b3:probability:independence}{独立}）。
求和：$\E T_n = \sum_k\frac n{n-k+1} = nH_n \sim n\ln n$；
$\V(T_n) = \sum\frac{1 - p_k}{p_k^2} \leq
n^2\sum_{j=1}^n\frac1{j^2} \leq \frac{\pi^2}6n^2$。

(b) 由切比雪夫不等式：$\P\bigl(\abs{T_n - nH_n} \geq \varepsilon
n\ln n\bigr) \leq \frac{\pi^2n^2/6}{\varepsilon^2n^2\ln^2n}
\to 0$，且 $\frac{nH_n}{n\ln n} \to 1$：$\frac{T_n}{n\ln n}
\to 1$ 依概率收敛。

(c) 由并界：$T_n > t$ 表示至少有一种类型在
$\lceil t\rceil$ 次抽取后仍未出现，故 $\P(T_n > t) \leq n(1 -
\frac1n)^{t} \leq n\,\eu^{-t/n}$；取 $t = \beta n\ln n$ 时：
$\leq n^{1 - \beta}$。当 $\beta > 1$ 时，$\sum_m
2^{m(1-\beta)} < \infty$，由博雷尔–坎泰利引理，沿 $n =
2^m$，几乎必然 $T_n \leq \beta n\ln n$ 最终成立——
特别地，对题中陈述的每个 $\beta > 2$ 都成立（任一 $\beta > 1$
沿该子列即可成立）。
\end{solution}

\begin{solution}{exo:b3:probability:10}
(a) 偶数位数字 $(b_{2k})_k$ 为\omterm{def:b3:probability:independence}{独立}同分布的公平
比特（\omterm{def:b3:probability:independence}{独立}数字族的一个子族），故 $U =
\sum_kb_{2k}2^{-k}$ 使每个二进区间获得正确的
概率（参见
\cref{thm:b3:probability:existence}）：均匀分布；$V$ 同理；
且 $(U, V)$ 依赖于不相交的数字块，故相互\omterm{def:b3:probability:independence}{独立}
（先在二进矩形上验证乘积，再由邓金定理）。

(b) $\Phi(\omega) = (U(\omega), V(\omega))$ 是可测的
且 $\Phi_*\lambda = \lambda\otimes\lambda = \lambda_2$
（两者在二进矩形上一致，由唯一性得）。交错
数字给出一个逆映射，它在去掉（零测）集后有定义，该零测集为
任一因子中的二进有理数的集合：这是一个保测
双射，介于 $\intcc01$ 的满\omterm{def:b3:measure:measure}{测度}子集与
$\intcc01^2$ 的满\omterm{def:b3:measure:measure}{测度}子集之间。与皮亚诺
（\cref{pb:b3:topology:1}）：连续性迫使满射，
却不保证单射；去掉连续性而仅求可测，
便换得\omterm{def:b3:measure:measure}{测度}同构——维数在\omterm{def:b3:measure:measure}{测度}论中不可见，
在\omterm{def:b3:topology:topology}{拓扑}中可见。
\end{solution}

\begin{solution}{exo:b3:probability:11}
(a) 分布的连续性使得并列成为零概率事件
（如本章次序统计量的论证所示），且
$n!$ 种 $(X_1, \dots, X_n)$ 的相对次序
是可交换的，故等可能。$R_n = 1$ 表示
最大值位于最后一个位置，概率为
$\frac{(n-1)!}{n!} = \frac1n$。

(b) 固定 $n$，以 $X_1,
\dots, X_{n-1}$ 的相对次序为条件：将 $X_n$ 插入 $n$ 个可能的
名次位置中是均匀的，且与该次序\omterm{def:b3:probability:independence}{独立}
（由 $n$ 元组的可交换性）。于是 $R_n$（事件
“$X_n$ 取得最高名次”）\omterm{def:b3:probability:independence}{独立}于整个
记录历史 $(R_1, \dots, R_{n-1})$，它是
前 $n - 1$ 个变量相对次序的函数。
由归纳法得全体\omterm{def:b3:probability:independence}{独立}，且 $\P(R_n = 1) =
\frac1n$。

(c) $\sum\P(R_n = 1) = \sum\frac1n = \infty$，且事件
相互\omterm{def:b3:probability:independence}{独立}，由博雷尔–坎泰利第二部分，记录
几乎必然地出现无穷多次（记录永不停止——但其以
对数速度变稀：$\E[\#\text{记录数} \leq n] =
H_n \approx \ln n$）。相邻记录：$\P(R_n = R_{n+1}
= 1) = \frac1{n(n+1)}$（由\omterm{def:b3:probability:independence}{独立}性），且
\[
\sum_n\frac1{n(n+1)} = \sum_n\Bigl(\frac1n -
\frac1{n+1}\Bigr) = 1 < \infty :
\]
应用博雷尔–坎泰利第一部分——仅有有限多个
相邻记录对出现，几乎必然。
\end{solution}

\begin{solution}{exo:b3:probability:12}
(a) 长为 $\ell$、始于位置 $i \leq n$ 的游程
其概率为 $2^{-\ell}$；由并界，$\P(L_n \geq \ell)
\leq n2^{-\ell}$。取 $\ell_n = (1 +
\varepsilon)\log_2n$ 时：$\P(L_n \geq \ell_n) \leq
n^{-\varepsilon}$。沿 $n = 2^k$：
$\sum_k2^{-k\varepsilon} < \infty$，故几乎必然 $L_{2^k} <
(1+\varepsilon)k$ 最终成立（博雷尔--坎泰利）；对一般的
$n$，取 $2^{k-1} < n \leq 2^k$，并利用
$L_n$ 的单调性及 $\log_22^{k-1} \leq \log_2n$：$L_n \leq L_{2^k}
< (1 + \varepsilon)k \leq (1 + \varepsilon)\frac{k}{k-1}
\log_2n$，额外因子可通过略微增大
$\varepsilon$ 来吸收。

(b) 取 $\ell = \lceil(1 - \varepsilon)\log_2n\rceil$ 及
$m = \lfloor n/\ell\rfloor$ 个互不相交的块：各块
相互\omterm{def:b3:probability:independence}{独立}，每块全为正面的概率为 $2^{-\ell}
\geq n^{-(1-\varepsilon)}/2$，故
\[
\P(L_n < \ell) \leq \bigl(1 - 2^{-\ell}\bigr)^{m}
\leq \exp\bigl(-m2^{-\ell}\bigr)
\leq \exp\Bigl(-c\,\frac{n^{\varepsilon}}{\log_2n}\Bigr)
\]
对某个常数 $c > 0$ 及充分大的 $n$，这些概率
沿 $n = 2^k$ 可和（事实上沿所有 $n$ 亦然）：
由博雷尔--坎泰利，几乎必然 $L_n \geq (1 -
\varepsilon)\log_2n$ 最终成立（单调性可在
相邻的 $2^k$ 之间补齐，如 (a) 所示，无碍）。

(c) 沿数列 $\varepsilon = \frac1j$ 两个界皆成立，
取可数多个全\omterm{def:b3:measure:measure}{测度}事件之交：
$\frac{L_n}{\log_2n} \to 1$ 几乎必然。当 $n = 10^6$ 时，$\log_2n
\approx 19.9$——长为 $\approx 20$ 的正面游程并非
可疑的异常，而是数学上的必然事件；若其缺席，
则是人为伪造“随机性”的证据（人类
极少敢连续写出超过 $5$ 或 $6$ 个正面）。
\end{solution}

\begin{solution}{pb:b3:probability:1}
\textbf{1.} $X_n^{\pm}$ 是 $X_n$ 的博雷尔函数：它们
仍两两\omterm{def:b3:probability:independence}{独立}
（见\cref{exo:b3:probability:3}(a)），同分布且\omterm{def:b3:lebesgue:l1}{可积}，
并有 $\E X_1 = \E X_1^+ - \E
X_1^-$。若定理对非负变量成立，
则对两部分分别应用并相减：
$\frac{S_n}n = \frac{S_n^+}n - \frac{S_n^-}n \to \E X_1^+ -
\E X_1^- = m$ 几乎必然。

\textbf{2.} $\P(X_n \neq Y_n) = \P(X_n > n) = \P(X_1 > n)$
（同分布），且 $\sum_n\P(X_1 > n) \leq \sum_n\P(X_1
\geq n) \leq \E X_1 < \infty$
（见\cref{exo:b3:product:3}(a)）。博雷尔--坎泰利（1）：几乎必然
$X_n = Y_n$ 对充分大的 $n$ 成立，因此 $S_n - S_n^*$
最终关于 $n$ 为常数：$\frac{S_n - S_n^*}n \to 0$
几乎必然，两个规范化和具有相同的
渐近行为。

\textbf{3.} $X_1\mathbf 1_{X_1 \leq n} \nearrow X_1$：由单调收敛定理
得 $\E Y_n \to m$；收敛数列的切萨罗平均
收敛于同一极限：$\frac{\E S_n^*}n =
\frac1n\sum_{k\leq n}\E Y_k \to m$。于是只需
证明 $\frac{S^*_n - \E S^*_n}{n} \to 0$ 几乎必然。

\textbf{4.} $\V(Y_n) \leq \E Y_n^2 = \E[X_1^2\mathbf
1_{X_1\leq n}]$。由级数形式的托内利定理，
\[
\sum_n\frac{\E[X_1^2\mathbf 1_{X_1\leq n}]}{n^2}
= \E\Bigl[X_1^2\!\!\sum_{n \geq \max(X_1, 1)}\!\frac1{n^2}
\Bigr]
\leq \E\Bigl[X_1^2\cdot\frac{4}{\max(X_1,1)}\Bigr]
\leq 4\,\E[X_1] < \infty,
\]
利用 $\sum_{n\geq x}n^{-2} \leq \frac4x$，其中 $x \geq 1$
（当 $x \geq 2$ 时：$\leq \frac1{x-1} \leq \frac2x$；当 $1
\leq x < 2$ 时：$\leq \frac{\pi^2}6 \leq \frac4x$，因为
$\frac4x > 2$），而 $X_1^2/\max(X_1, 1) \leq X_1$ 在
$X_1 \gtrless 1$ 两种情形下均成立。

\textbf{5.} 两两\omterm{def:b3:probability:independence}{独立}给出 $\E[(Y_i - \E
Y_i)(Y_j - \E Y_j)] = 0$（$i \neq j$），此为两变量的乘积公式，
故方差可加：$\V(S^*_k) =
\sum_{n\leq k}\V(Y_n)$。对每个 $k_j$ 用切比雪夫不等式，
再求和：
\[
\sum_j\P\Bigl(\abs{S^*_{k_j} - \E S^*_{k_j}} \geq
\varepsilon k_j\Bigr)
\leq \frac1{\varepsilon^2}\sum_j\frac1{k_j^2}\sum_{n\leq
k_j}\V(Y_n)
= \frac1{\varepsilon^2}\sum_n\V(Y_n)\!\!\sum_{j : k_j\geq
n}\!\frac1{k_j^2}
\]
（对非负二重级数用托内利定理）。

\textbf{6.} $k_j = \lfloor\alpha^j\rfloor \geq
\frac{\alpha^j}2$（只要 $\alpha^j \geq 1$，即对所有
$j \geq 0$ 成立，$\lfloor x\rfloor \geq \frac x2$ 当 $x \geq
1$）。由此
\[
\sum_{j : k_j \geq n}\frac1{k_j^2}
\leq 4\sum_{j : \alpha^j \geq n}\alpha^{-2j}
\leq \frac{4}{1 - \alpha^{-2}}\cdot\frac1{n^2}
= \frac{C_\alpha}{n^2},
\]
（从第一个 $j$ 满足 $\alpha^j \geq
n$ 起的几何级数）。结合第 4--5 题，该二重和
有限；博雷尔--坎泰利（1）对每个有理数
$\varepsilon$ 应用并取交，得
$\frac{S^*_{k_j} - \E S^*_{k_j}}{k_j} \to 0$ 几乎必然，并结合
第 3 题：$\frac{S^*_{k_j}}{k_j} \to m$ 几乎必然。

\textbf{7.} 由 $Y_n \geq 0$ 知 $n \mapsto S^*_n$
非减：对 $k_j \leq n \leq k_{j+1}$，
\[
\frac{S^*_{k_j}}{k_{j+1}} \leq \frac{S^*_n}{n} \leq
\frac{S^*_{k_{j+1}}}{k_j},
\]
此即在两边分别乘以
$\frac{k_j}{k_{j+1}}$ 与 $\frac{k_{j+1}}{k_j}$ 后所得的夹逼式。由于
$\frac{k_{j+1}}{k_j} \to \alpha$，由第 6 题几乎必然有
\[
\frac m\alpha \leq \liminf_n\frac{S^*_n}n \leq
\limsup_n\frac{S^*_n}n \leq \alpha m .
\]

\textbf{8.} 对 $\alpha = 1 + \frac1p$ 应用第 7 题，$p
\in \N^*$：可数多个几乎必然事件；
在它们的交上令 $p \to \infty$：$\lim\frac{S^*_n}n =
m$ 几乎必然。结合第 1--3 题，$\frac{S_n}n \to \E X_1$
几乎必然：这是两两\omterm{def:b3:probability:independence}{独立}条件下的
强大数律。

\textbf{9.} \omterm{def:b3:probability:independence}{独立}型假设共出现
三次：(i) 方差的可加性（第 5 题）——
两两\omterm{def:b3:probability:independence}{独立}即可；(ii) 同分布，出现在
截断和（第 2 题）与均值计算
（第 3 题）——完全不需要任何\omterm{def:b3:probability:independence}{独立}性；(iii)
博雷尔--坎泰利（1）（第 2 题与第 6 题）——无需
任何\omterm{def:b3:probability:independence}{独立}性即成立。完全互相\omterm{def:b3:probability:independence}{独立}从未被
用到：这是埃特马迪的观察。

\textbf{10.} 固定底数 $b$ 与数字 $r$。$b$ 进制
数字 $(d_k)$ 取自均匀的 $\omega$，\omterm{def:b3:probability:independence}{独立}同分布，
在 $\{0, \dots, b-1\}$ 上服从均匀分布（数字向量的每个取值占据一个
长度为 $b^{-m}$ 的区间：其论证与
\cref{thm:b3:probability:existence} 原样相同）。将强大数律用于\omterm{def:b3:probability:independence}{独立}同分布的
有界变量 $\mathbf
1_{d_k = r}$，结论为：几乎必然，数字 $r$ 的频率
趋于 $\frac1b$。对可数多个
数对 $(b, r)$ 取交：几乎每个数都是
\emph{在每一进制下都单正规}。一个显式的非正规数为
$x = 0.100100100\ldots_2$（数字 1 的频率
$\frac13 \neq \frac12$）。这一对比发人深省：几乎
所有数都是正规的，然而对于 $\sqrt2$、$\eu$ 或 $\pi$，
正规性仍未得证——\omterm{def:b3:measure:measure}{测度}论只能计数，
而不能具体给出。

\textbf{11.} 由 \cref{exo:b3:probability:10} 反复应用，单个
均匀变量可产生一列\omterm{def:b3:probability:independence}{独立}同分布的
均匀\emph{向量} $U_k$，位于 $\intcc01^d$ 上（
每个 $U_n$ 的数字集出自
\cref{thm:b3:probability:existence}，拆为 $d$ 个子族）。
对 $g \in L^1(\intcc01^d)$，变量 $g(U_k)$
是\omterm{def:b3:probability:independence}{独立}同分布的\omterm{def:b3:lebesgue:l1}{可积}变量，均值为 $\int g\,\dd\lambda_d$
（由变量变换）：强大数律给出
\[
\frac1n\sum_{k=1}^ng(U_k)
\xrightarrow[n\to\infty]{\text{几乎必然}}
\int_{\intcc01^d}g\,\dd\lambda_d :
\]
\omterm{ex:b3:probability:sllnapps}{蒙特卡罗}积分几乎必然收敛，且在任意
维数下均成立——误差大小属于中心极限定理的
范畴（见\cref{ch:b3:clt}）。

\textbf{12.} 记 $A_k = \{\abs{S_k} \geq \varepsilon\} \cap
\bigcap_{j<k}\{\abs{S_j} < \varepsilon\}$：$A_k$ 互不
相交，且其并为 $A = \{\max_{k\leq n}\abs{S_k} \geq
\varepsilon\}$。则
\[
\E S_n^2 \geq \sum_{k=1}^n\E\bigl[S_n^2\mathbf 1_{A_k}\bigr]
= \sum_{k=1}^n\E\Bigl[\bigl(S_k^2 + 2S_k(S_n - S_k) + (S_n
- S_k)^2\bigr)\mathbf 1_{A_k}\Bigr]
\geq \sum_{k=1}^n\E\bigl[S_k^2\mathbf 1_{A_k}\bigr],
\]
因交叉项为零，$S_k\mathbf 1_{A_k}$ 是
关于组 $(Z_1, \dots, Z_k)$ 的博雷尔函数，它与
$S_n - S_k$ \omterm{def:b3:probability:independence}{独立}，后者是 $(Z_{k+1},
\dots, Z_n)$ 的函数（见\cref{thm:b3:probability:independence}），因此
$\E[S_k\mathbf 1_{A_k}(S_n - S_k)] = \E[S_k\mathbf
1_{A_k}]\,\E[S_n - S_k] = 0$。在 $A_k$ 上，$S_k^2 \geq
\varepsilon^2$，从而 $\E S_n^2 \geq
\varepsilon^2\sum_k\P(A_k) = \varepsilon^2\P(A)$；而 $\E
S_n^2 = \sum_{k\leq n}\V(Z_k)$（方差可加）。
决定性的一步是这一分解：$S_k\mathbf 1_{A_k}$ 是
关于整个第一组的一个\emph{非线性}函数，而它与第二组
的\omterm{def:b3:probability:independence}{独立}性即为组\omterm{def:b3:probability:independence}{独立}性，而
仅有 $Z_i$ 两两\omterm{def:b3:probability:independence}{独立}——只能
使成对去相关，不足以证明此点。

\textbf{13.} 固定 $N$，对 $Z_{N+1},
\dots, Z_{N+m}$ 应用第 12 题：
\[
\P\Bigl(\max_{N < k \leq N+m}\abs{S_k - S_N} >
\varepsilon\Bigr) \leq
\frac1{\varepsilon^2}\sum_{j=N+1}^{N+m}\V(Z_j) \leq
\frac{r_N}{\varepsilon^2},
\qquad r_N = \sum_{j>N}\V(Z_j) .
\]
事件随 $m$ 递增；由下连续性得
$\P(\sup_{k>N}\abs{S_k - S_N} > \varepsilon) \leq
r_N/\varepsilon^2$，且由假设 $r_N \to 0$。因此
对每个 $p \in \N^*$，$\P\bigl(\bigcap_N\{\sup_{k>N}
\abs{S_k - S_N} > \frac1p\}\bigr) \leq \inf_Np^2r_N = 0$：
几乎必然地，对每个 $p$ 存在 $N$，使
$\sup_{k>N}\abs{S_k - S_N} \leq \frac1p$（对
关于 $p$ 的可数多个几乎必然事件取交），从而 $\abs{S_k -
S_l} \leq \frac2p$ 对一切 $k, l > N$ 成立：部分和
几乎必然构成柯西列，从而几乎必然收敛。

\textbf{14.} 变量 $Z_n = x_n\varepsilon_n$ 是
\omterm{def:b3:probability:independence}{独立}的（\omterm{def:b3:probability:independence}{独立}变量的博雷尔函数，
见\cref{exo:b3:probability:3}(a)），中心化的，且 $\V(Z_n) =
x_n^2$：第 13 题适用于 $\sum_nx_n^2 < \infty$ 的情形，
并给出几乎必然收敛。一般地，对每个 $N$，
$\sum_nx_n\varepsilon_n$ 的收敛性不受
$\varepsilon_1, \dots, \varepsilon_N$ 取值的影响：
收敛事件属于尾 $\sigma$-代数，
即\omterm{def:b3:probability:independence}{独立}序列 $(\varepsilon_n)$ 的尾，故由柯尔莫哥洛夫
零一律（见\cref{thm:b3:probability:zeroone}）迫使
其概率为 $0$ 或 $1$。

\textbf{15.} (a) 在水平 $\theta\E Z$ 处分割，并
对上半部分使用柯西--施瓦茨不等式，
\[
\E Z = \E\bigl[Z\mathbf 1_{Z \leq \theta\E Z}\bigr] +
\E\bigl[Z\mathbf 1_{Z > \theta\E Z}\bigr]
\leq \theta\,\E Z + \sqrt{\E Z^2}\,
\sqrt{\P(Z > \theta\E Z)} ,
\]
于是 $(1 - \theta)\E Z \leq \sqrt{\E Z^2\,\P(Z > \theta\E
Z)}$；两边平方。(b) 展开 $T_n^4 =
\sum_{i,j,k,l}x_ix_jx_kx_l\,
\E[\varepsilon_i\varepsilon_j\varepsilon_k\varepsilon_l]$:
\omterm{def:b3:probability:space}{期望}为 $1$，当指标两两配对时（四个指标全
相等，或两对不同指标，后者有 $3$
种排列），否则\omterm{def:b3:probability:space}{期望}为 $0$（未配对的符号均值为零，
且由\omterm{def:b3:probability:independence}{独立}性可提出因子）。于是
\[
\E T_n^4 = \sum_kx_k^4 + 3\sum_{i\neq j}x_i^2x_j^2 =
3s_n^4 - 2\sum_kx_k^4 \leq 3s_n^4 .
\]
(c) 由佩利--齐格蒙德不等式，取 $Z = T_n^2$，$\E Z = s_n^2$，
$\theta = \frac14$:
\[
\P\Bigl(\abs{T_n} > \frac{s_n}2\Bigr) = \P\Bigl(T_n^2 >
\frac{s_n^2}4\Bigr) \geq \Bigl(\frac34\Bigr)^2
\frac{s_n^4}{3s_n^4} = \frac3{16} .
\]
若该级数以正概率收敛，则它
几乎必然收敛（第 14 题），因此 $\sup_n\abs{T_n} <
\infty$ 几乎必然，且存在某个 $M$ 满足
$\P(\sup_n\abs{T_n} > M) < \frac3{16}$；但只要 $s_n
> 2M$, $\P(\abs{T_n} > M) \geq \P(\abs{T_n} > \frac{s_n}2)
\geq \frac3{16}$：矛盾。因此发散几乎
必然；结合第 14 题，二分性即告完整。

\textbf{16.} 此处 $x_n = n^{-s}$，且 $\sum_nn^{-2s} <
\infty$ 当且仅当 $s > \frac12$ 时成立：由第 14--15 题，
$\sum_n\frac{\varepsilon_n}{n^s}$ 几乎必然收敛，当且
仅当 $s > \frac12$（对 $s \leq \frac12$，几乎必然
发散）。对 $\frac12 < s \leq 1$，收敛
绝不是绝对收敛的。这一比较颇有启发：完全交替的符号
以强度 $n^{-s}$ 抵消，对每个 $s
> 0$，而典型的随机符号仅以
平方根强度——第 21 题的随机游走
增长为 $\sqrt n$，阿贝尔求和法恰好将这种增长
转化为 $\sum\varepsilon_nn^{-s}$ 的收敛性，
其中 $s > \frac12$。

\textbf{17.} (a) $\cosh\lambda =
\sum_k\frac{\lambda^{2k}}{(2k)!}$ 与 $\eu^{\lambda^2/2} =
\sum_k\frac{\lambda^{2k}}{2^kk!}$；且 $(2k)! \geq 2^kk!$
逐项成立，因为 $\frac{(2k)!}{k!} =
\prod_{i=1}^k(k + i) \geq \prod_{i=1}^k(2i) = 2^kk!$（每个
因子满足 $k + i \geq 2i$（$i \leq k$），从而
事实上 $(2k)! \geq 2^k(k!)^2 \geq 2^kk!$。(b) 注意 $a \leq 0 \leq b$（$Z$ 已
中心化），由 $z \mapsto \eu^{\lambda z}$ 的凸性，
对 $z \in \intcc ab$：
\[
\eu^{\lambda z} \leq \frac{b - z}{b - a}\,\eu^{\lambda a} +
\frac{z - a}{b - a}\,\eu^{\lambda b},
\qquad\text{于是}\qquad
\E\,\eu^{\lambda Z} \leq \frac{b\,\eu^{\lambda a} -
a\,\eu^{\lambda b}}{b - a}
= (1 - p)\eu^{-pt} + p\,\eu^{(1-p)t} = \eu^{\varphi(t)}
\]
其中 $p = \frac{-a}{b-a} \in \intcc01$，$t = \lambda(b -
a)$，$\varphi(t) = -pt + \log(1 - p + p\eu^t)$。则
$\varphi(0) = 0$, $\varphi'(t) = -p + \frac{p\eu^t}{1 - p +
p\eu^t}$ 在 $0$ 处为零，且 $\varphi''(t) = \rho(1 -
\rho) \leq \frac14$，其中 $\rho = \frac{p\eu^t}{1 - p +
p\eu^t} \in \intcc01$：由 $2$ 阶泰勒公式得
$\varphi(t) \leq \frac{t^2}8 = \frac{\lambda^2(b-a)^2}8$.

\textbf{18.} 当 $\lambda > 0$ 时，对正变量
$\eu^{\lambda(S_n - \E S_n)}$ 应用马尔可夫不等式
（\cref{prop:b3:probability:markov}），并由\omterm{def:b3:probability:independence}{独立}变量的
乘积公式得
\[
\P(S_n - \E S_n \geq t) \leq \eu^{-\lambda
t}\prod_{i=1}^n\E\,\eu^{\lambda(X_i - \E X_i)}
\leq \exp\Bigl(-\lambda t +
\frac{\lambda^2}8\sum_i(b_i - a_i)^2\Bigr),
\]
由第 17(b) 题，对每个中心化变量 $X_i - \E X_i
\in \intcc{a_i - \E X_i}{b_i - \E X_i}$（宽度相同）。
令指数在 $\lambda = \frac{4t}{D}$ 处取最小值，其中 $D =
\sum_i(b_i - a_i)^2$，得 $-\frac{2t^2}D$。下
尾可对 $(-X_i)$ 应用上述结论得到。

\textbf{19.} 取 $t = n\varepsilon$ 与 $D = n(b - a)^2$：
\[
\P\Bigl(\Bigl|\frac{S_n}n - m\Bigr| \geq \varepsilon\Bigr)
\leq 2\exp\Bigl(\frac{-2n^2\varepsilon^2}{n(b-a)^2}\Bigr)
= 2\exp\Bigl(\frac{-2n\varepsilon^2}{(b-a)^2}\Bigr),
\]
该界关于 $n$ 可和（类几何级数）。
由博雷尔--坎泰利（\cref{thm:b3:probability:borelcantelli}）知，
几乎必然有 $\abs{\frac{S_n}n - m} < \varepsilon$
最终成立；对 $\varepsilon = \frac1p$ 求交，
便得 $\frac{S_n}n \to m$ 几乎必然。比较：埃特马迪仅要求
$X_1 \in L^1$ 与两两\omterm{def:b3:probability:independence}{独立}，且不给出速率；霍夫丁
则要求有界与完全\omterm{def:b3:probability:independence}{独立}，并在每个
有限 $n$ 处给出显式的指数型保证——
两个定理回答的是
同一极限的不同问题。

\textbf{20.} $g(U_k)$ 为\omterm{def:b3:probability:independence}{独立}同分布的变量，取值于
$\intcc01$ 且均值为 $\int g\,\dd\lambda_d$（由变量替换），故
由第 18 题，取 $b_i - a_i = 1$，$t = n\varepsilon$，给出
双侧界 $2\eu^{-2n\varepsilon^2} \leq \delta$，
只要 $\eu^{2n\varepsilon^2} \geq \frac2\delta$，即
$n \geq \frac{\log(2/\delta)}{2\varepsilon^2}$。对于
$\varepsilon = \delta = 10^{-2}$:
\[
n \geq \frac{\log 200}{2\cdot10^{-4}} =
\frac{5.2983\ldots}{0.0002} \approx 26\,492 :
\]
约 $26\,500$ 个样本即可保证 $1\%$ 的精度，
并具有 $99\%$ 的置信度——在每个维数 $d$ 下，对每个
取值于 $\intcc01$ 的可测被积函数均成立。
第 11 题的强大数定律保证了收敛，但没有有限 $n$ 的
保证；确定性网格（每个坐标轴取 $k$ 个点）
需要 $k^d$ 次求值，关于 $d$ 呈指数增长。集中现象
使\omterm{ex:b3:probability:sllnapps}{蒙特卡罗}成为\emph{方法}，而不只是
一种希望。

\textbf{21.} \omterm{def:b3:probability:independence}{独立}性与乘积公式：
$\E\,\eu^{\lambda S_n} = (\E\,\eu^{\lambda\varepsilon_1})^n
= (\cosh\lambda)^n \leq \eu^{n\lambda^2/2}$ 由第
17(a) 题，对 $\eu^{\lambda S_n}$ 用马尔可夫不等式：
\[
\P(S_n \geq x) \leq \eu^{-\lambda x + n\lambda^2/2}
= \eu^{-x^2/(2n)}
\qquad\text{ 在最优 } \lambda = \frac xn,
\]
对 $-S_n$（分布相同）的对称界同样成立，使
$\abs{S_n}$ 的常数加倍。

\textbf{22.} 固定 $\eta > 0$，并令 $x_n = (1 +
\eta)\sqrt{2n\log n}$，其中 $n \geq 2$：
\[
\P(\abs{S_n} \geq x_n) \leq 2\exp\bigl(-(1 +
\eta)^2\log n\bigr) = \frac{2}{n^{(1+\eta)^2}},
\]
级数因 $(1 + \eta)^2 > 1$ 而收敛。由博雷尔–坎特利引理，几乎必然
$\abs{S_n} < (1 + \eta)\sqrt{2n\log n}$ 对所有充分大的 $n$ 成立，
故 $\limsup_n\frac{\abs{S_n}}{\sqrt{2n\log n}} \leq 1 +
\eta$ 几乎必然；对 $\eta =
\frac1p$，$p \in \N^*$，取这些几乎必然事件的交即得结论。规模为
$n$ 的游走典型振幅为 $\sqrt n$（即其方差），而
即便其最坏的偏离，也至多超出该尺度的
$\sqrt{2\log n}$ 倍。

\textbf{23.} 取 $n_j = 2^j$ 及 $x = (1 +
\eta)\sqrt{2n_j\log\log n_j}$（$j \geq 2$ 时有定义），由第
21 题可得
\[
\P\bigl(S_{n_j} \geq x\bigr) \leq \exp\bigl(-(1 +
\eta)^2\log\log n_j\bigr) = (j\log 2)^{-(1+\eta)^2},
\]
关于 $j$ 收敛，因 $(1 + \eta)^2 > 1$。由博雷尔–坎特利
引理及 $\eta = \frac1p$ 给出 $\limsup_jS_{n_j}/\sqrt{2n_j
\log\log n_j} \leq 1$ 几乎必然。完整上半部分所缺的是检查点之间的
衔接：须
证明 $\max_{n_j\leq n\leq n_{j+1}}S_n$ 超过
$(1+\eta)\sqrt{2n_j\log\log n_j}$ 的情形仅有限多次，
这需要一个具有\emph{高斯}型尾的极大不等式
（列维反射不等式，或奥塔维亚尼
不等式，此处不予证明）。第 12 题在数量上
过弱，它给出的概率上界为
\[
\frac{n_j}{(1+\eta)^2\,2n_j\log\log n_j}
= \frac{1}{2(1+\eta)^2\log(j\log2)},
\]
该界趋于 $0$，但关于 $j$ \emph{发散}：
故博雷尔–坎特利引理无法得出结论。重对数律的下半部分
是把第二
博雷尔–坎特利引理用于\omterm{def:b3:probability:independence}{独立}增量
$S_{n_{j+1}} - S_{n_j}$，并利用
高斯型尾的匹配下界。两项精细化都是真正的本科三年级概率论，
比本题高一门课程；本题\omterm{def:b3:probability:independence}{独立}完成的是
精确的重对数尺度，
沿几何时刻而言。

\textbf{24.} 每个 $\hat p_i$ 是 $n$ 个\omterm{def:b3:probability:independence}{独立}同分布
的示性变量的平均，这些变量取值于 $\intcc01$，\omterm{def:b3:probability:space}{期望}为
$\P(A_i)$。由霍夫丁不等式，$\P(\abs{\hat p_i - \P(A_i)} >
\varepsilon) \leq 2\eu^{-2n\varepsilon^2}$。取并界，
乘以 $N$。解 $2N\eu^{-2n\varepsilon^2} \leq
\delta$：$n \geq \frac{\ln(2N/\delta)}{2\varepsilon^2}$。
数值上：$\ln\frac{2\cdot10^6}{0.05} =
\ln(4\cdot10^7) \approx 17.5$，故 $n \geq
\frac{17.5}{2\cdot10^{-4}} \approx 87\,600$：估计
\emph{一个}概率到 $\pm1\%$ 的精度约需 $18\,500$
个样本（$\ln(2/\delta)/2\varepsilon^2$），而
\emph{一百万}个概率只需多 $\approx 4.7$ 倍的样本，
——一致性的代价是 $\ln N$ 而非 $N$，正是这一观察
使经验风险最小化，并由此使
机器学习在统计上成为可能。

\textbf{25.} 变量 $X_n = \frac{\varepsilon_n}
{n^\alpha}$ \omterm{def:b3:probability:independence}{独立}、中心化、有界，
且 $\sum_n\V(X_n) = \sum_nn^{-2\alpha}$。若 $\alpha >
\frac12$：方差级数收敛，而
由单级数定理（第六部分）给出
$\sum X_n$ 几乎必然收敛。若 $\alpha \leq \frac12$：方差级数
发散，逆方向的一半（第六部分的
佩利–齐格蒙德论证，因求和项
以 $1$ 为界而适用）给出几乎必然发散。绝对收敛要求
$\sum n^{-\alpha} < \infty$，即 $\alpha > 1$。
在 $\intoc{\frac12}1$ 上，级数几乎必然收敛，
尽管 $\sum\abs{X_n} = \infty$ 恒成立：各项的符号
相互配合，几乎必然抵消——这种靠抵消实现的
收敛，任何绝对收敛判别法都察觉不到，并且
（由零一律）结论依然是确定的。
\end{solution}
